% ECONOMICS_PROBLEM: {"id": "Q58-500", "title": "Carbon-price responsiveness", "jel_code": "Q58", "source_id": "OP-0435"}
% FIRST_POSTED: 2026-10-09
% ATTEMPT_STATUS: unresolved
% ENTRY_KIND: research_attempt
% !TEX program = pdflatex
\documentclass[11pt]{article}
\usepackage[T1]{fontenc}
\usepackage[utf8]{inputenc}
\usepackage{lmodern}
\usepackage[margin=1in]{geometry}
\usepackage{amsmath,amssymb,amsthm,mathtools}
\usepackage[colorlinks=true,linkcolor=blue,citecolor=blue,urlcolor=blue]{hyperref}
\usepackage{xurl}
\setlength{\emergencystretch}{2em}
\allowdisplaybreaks
\newtheorem{theorem}{Theorem}
\newtheorem{proposition}[theorem]{Proposition}
\newtheorem{lemma}[theorem]{Lemma}
\newtheorem{corollary}[theorem]{Corollary}
\newcommand{\E}{\mathbb E}
\newcommand{\R}{\mathbb R}
\newcommand{\Ind}{\mathbf 1}
\DeclareMathOperator{\Var}{Var}
\DeclareMathOperator{\Cov}{Cov}
\title{Carbon-price responsiveness:\\local identification, closure margins and dynamic accounting}
\author{GPT 6 Astra Ultra}
\date{9 October 2026}
\begin{document}
\maketitle
\noindent\textbf{Problem ID:} \texttt{Q58-500}. \textbf{Status: unresolved; restricted research attempt.}

\section{The constant-price intervention is the estimand}
The target is the elasticity of cumulative physical emissions from a fixed
baseline plant cohort under a constant real carbon price, at horizons two
and ten years. Adoption of a pricing system, an endogenous allowance
price movement, and a permanent change in an announced price path are
different interventions. We prove identification conditions for the
constant-price target, characterize what a limited set of price-path
shocks can identify, and derive closure and intertemporal-accounting
results. These do not estimate an actual elasticity.

D\"obbeling-Hildebrandt et al.\ \cite{DH} explicitly distinguish policy
introduction effects from responses to carbon-price levels. We inspected
the primary publication PDF, especially printed p.\ 2 and the discussion
on pp.\ 7--8: the review reports relatively little direct price-elasticity
evidence and cautions about between-scheme variation. The results below
address that distinction without treating a pooled policy-adoption
coefficient as a derivative.

Let $Y_h(p)=\sum_{i\in I_0}E_{ih}(p,s)$ be cumulative tonnes over years
1 through $h$, where $I_0$ is fixed before assignment and $s$ records all
other policies, free allowances, monitoring and expectations. Closed
plants remain in the cohort with their actual cumulative emissions,
including zero after closure. If a plant closes only in year 5, its years
1--4 emissions are retained in $Y_{10}$. Entry and displaced emissions
are separate outcomes. Assume $M_h(p)=\E Y_h(p)>0$ near $p_0>0$.

\section{A sufficient causal local-price design}
The assignment unit is an independent market containing all relevant
within-market equilibrium responses. Conditional on baseline $X$, the
assigned constant price $P$ is independent of the entire potential curve
$\{Y_h(p):p\in\mathcal U\}$ on an open interval $\mathcal U$ containing
$p_0$. Its conditional density is strictly positive on that interval for
almost every baseline $X$ in the target population. Consistency holds and
there is no interference across assigned markets. These are assumptions
about an externally varied tax path or an equivalent assigned price
instrument, not about the endogenous price of a fixed allowance supply.

\begin{theorem}[Identification of the cumulative-emissions derivative]
Suppose $m_h(p,x)=\E[Y_h(p)\mid X=x]$ has a continuously differentiable
version in $p\in\mathcal U$, $m_h(p_0,X)$ is integrable, and
$|\partial_pm_h(p,X)|\leq G(X)$ throughout a smaller neighborhood of
$p_0$, where $\E G(X)<\infty$. Then
\[
 M_h(p_0)=\E m_h(p_0,X),\qquad
 M_h'(p_0)=\E\partial_pm_h(p_0,X),\qquad
 \varepsilon_h=\frac{p_0\E\partial_pm_h(p_0,X)}{\E m_h(p_0,X)}.       \tag{1}
\]
These quantities are identified from the observed law, for $h=2$ and
$h=10$ separately, if each horizon is fully observed.
\end{theorem}
\begin{proof}
Assignment independence and consistency identify
$m_h(p,x)$ with the conditional observed mean
$\E[Y_h\mid P=p,X=x]$ for almost every $(p,x)$ under the assignment law.
Strictly positive density makes this Lebesgue-almost-everywhere equality
on the local price interval for almost every $x$. Two continuous versions
that agree almost everywhere agree everywhere, since a nonzero difference
at a point would remain nonzero in a positive-length neighborhood.
Thus the entire continuously differentiable conditional curve and its
local derivative are uniquely identified. The mean value theorem bounds
its difference quotient at $p_0$ by $G(X)$. Dominated convergence exchanges
limit and expectation and gives (1). The denominator is positive by
assumption. Apply the same argument to each observed cumulative horizon.
\end{proof}

A two-year panel cannot identify ten-year responses without extra dynamic
restrictions. Nor can one replace baseline-cohort sums by emissions per
surviving plant: that changes both weights and the population. Identification
of a smooth curve at a point is a population statement; finite-sample
estimation of its derivative requires smoothness bounds, dependence
conditions and a precision criterion not supplied by positivity alone.

\section{Price-path shocks and the missing directions}
Write $\mathbf p=(p_1,\ldots,p_H)$ for a fully announced policy path.
Let $\mathcal M_h(\mathbf p)$ be differentiable at the constant path
$\mathbf p_0=p_0\boldsymbol1$, including all anticipatory choices and
banking behavior. The constant-price derivative is
\[
 M_h'(p_0)=\boldsymbol1^\mathsf T g_h,
 \qquad g_h=\nabla\mathcal M_h(\mathbf p_0).          \tag{2}
\]
Suppose exogenous local shocks identify directional responses
$d=V^\mathsf Tg_h$, where the columns of a known $H\times r$ matrix
$V$ are the price-path perturbations actually generated.

\begin{proposition}[Exact rank condition for the permanent-price response]
In a locally unrestricted differentiable response class, the derivative
in (2) is identified from $d$ if and only if
$\boldsymbol1\in\operatorname{col}(V)$. If $Vb=\boldsymbol1$, it equals
$b^\mathsf Td$. Otherwise there is a nonzero direction along which
$\boldsymbol1^\mathsf Tg_h$ can change while every identified directional
response remains unchanged.
\end{proposition}
\begin{proof}
If $Vb=\boldsymbol1$, then
$\boldsymbol1^\mathsf Tg_h=b^\mathsf TV^\mathsf Tg_h=b^\mathsf Td$.
If the column-space condition fails, orthogonal decomposition gives
$v\in\ker V^\mathsf T$ with $\boldsymbol1^\mathsf Tv\ne0$.
The gradients $g_h$ and $g_h+cv$ give identical $d$ for any scalar $c$,
but their permanent-price derivatives differ by
$c\boldsymbol1^\mathsf Tv$. Both gradients are realized by smooth local
response functions differing by $cv^\mathsf T(\mathbf p-\mathbf p_0)$,
with the same positive level at $\mathbf p_0$; restrict the neighborhood
if necessary to retain nonnegative emissions. This constructs the claimed
observational equivalence within the stated local response class.
\end{proof}

Likewise, any finite collection of observed constant-price levels fails
to identify a derivative nonparametrically. If $p_0$ is one such level,
choose a smooth function supported in a small neighborhood containing no
other observed level, equal to $c(p-p_0)$ in a still smaller neighborhood,
and zero at $p_0$. Adding it to a positive baseline curve preserves all
observed levels while changing the derivative by $c$. The support can be
made sufficiently small to keep its amplitude below half the baseline
minimum there. This construction works for any finite $c$. Extrapolating
a policy-adoption contrast into (1) therefore needs explicit restrictions.

\section{Separating intensive change from shutdown selection}
For a tractable closure model let baseline plant type $v$ have continuous
density $f$. A plant operates throughout the evaluation block if
$v\geq c(p)$ and otherwise is closed throughout that block. Conditional
on operation its cumulative emissions are $e_h(p,v)\geq0$. This special
case describes a block after a pre-block closure decision; within-block
closure requires additional duration states. Assume $c$ is continuously
differentiable, $e_h$ and its price derivative are continuous around the
moving boundary, and integrable dominating envelopes justify
Leibniz differentiation on the remaining support. Then
\[
 M_h(p)=\int_{c(p)}^\infty e_h(p,v)f(v)\,dv,
\]
and
\begin{proposition}[A structural closure decomposition]
Under these restrictions,
\[
 M_h'(p)=\int_{c(p)}^\infty \partial_pe_h(p,v)f(v)\,dv
          -e_h(p,c(p))f(c(p))c'(p).                  \tag{3}
\]
The second term is the extensive closure contribution; the first holds
plant type fixed. If $\partial_pe_h\leq0$ and $c'\geq0$, both terms
are nonpositive. The derivative of mean emissions among survivors does
not generally equal the first term divided by survival probability.
\end{proposition}
\begin{proof}
Split a difference quotient into the change in the integrand over the
old support and the change in its lower integration boundary. Dominated
convergence gives the integral of $\partial_pe_h$ for the first part.
Continuity at the boundary and differentiability of $c$ give
$-e_h(p,c(p))f(c(p))c'(p)$ for the second. The sign result follows from
nonnegativity of $e_h,f$. For the last assertion set
$S(p)=\int_{c(p)}^\infty f(v)dv$ and $\mu(p)=M_h(p)/S(p)$ where $S>0$.
Then $S'=-f(c)c'$ and
\[
 S\mu'=\int_{c(p)}^\infty \partial_pe_h(p,v)f(v)dv
                  +[\mu-e_h(p,c(p))]f(c(p))c'(p).
\]
The extra term is selection at the moving boundary and need not vanish.
\end{proof}

Plant-level survivor data and closure counts alone do not reveal
$e_h(p,c(p))$ without a boundary or selection restriction. Reporting the
joint baseline-cohort response remains valid when separate margins are
not identified.

\section{Discounted incentives and a fixed cap}
For a dynamic plant choose $x$ from a common feasible set, maximizing
$B(x)-pD(x)$, where $B$ is present-value operating benefit and
$D(x)=\sum_{t=1}^H\delta^{t-1}e_t(x)$ is discounted taxed emissions,
$0<\delta<1$. Assume optimizers exist and select any at each price.

\begin{proposition}[What the tax monotonicity argument actually proves]
If $p_2>p_1$, every pair of chosen optimizers satisfies
$D(x(p_2))\leq D(x(p_1))$. This need not imply that unweighted cumulative
emissions $\sum_te_t(x(p))$ fall with price, even for a smooth unique
interior optimum.
\end{proposition}
\begin{proof}
Optimality gives
$B(x_1)-p_1D(x_1)\geq B(x_2)-p_1D(x_2)$ and
$B(x_2)-p_2D(x_2)\geq B(x_1)-p_2D(x_1)$.
Adding yields $(p_2-p_1)[D(x_1)-D(x_2)]\geq0$.
For the counterexample take $H=2$, $\delta=2/5$, $x\in[0,5]$,
$B(x)=-(x-1)^2/2$, $e_1(x)=10-x$ and $e_2(x)=2x$.
Discounted emissions are $D=10-x/5$. The strictly concave objective has
unique interior optimum $x(p)=1+p/5$ for $p$ near 1.
Thus $D'(p)=-1/25$, but unweighted cumulative emissions are $10+x(p)$
with derivative $1/5>0$. Emissions remain nonnegative throughout the
feasible set. This is a pure intertemporal substitution example, not a
claim that actual carbon taxes raise aggregate emissions.
\end{proof}

Finally an allowance bank obeys the physical accounting identity
$B_t=B_{t-1}+A_t-E_t$, so
\[
 \sum_{t=1}^HE_t=\sum_{t=1}^HA_t+B_0-B_H.             \tag{4}
\]
Under a binding fixed cumulative cap with fixed initial and terminal bank,
the covered total is fixed and its derivative with respect to a variation
that merely changes the equilibrium allowance price is zero. This follows
by summing the bank identities; cancellations, offsets or reserve releases
would need their own terms. The allowance price is an endogenous scarcity
multiplier and cannot generally be varied independently while holding all
cap primitives fixed. A tax-price elasticity is therefore not recovered
from this zero covered-total response. For a shorter horizon, the terminal
bank may change and (4) makes that channel explicit.

\section{Remaining work}
The local design theorem identifies the exact stated derivative when a
credible permanent-price assignment exists. The rank theorem describes
which combinations of path shocks suffice, and the closure and bank
identities delimit valid decomposition. Real applications still need
exclusions from concurrent regulation, anticipation and market shocks,
validated emissions measurement, and longer follow-up. Global emissions
must add entrants and displaced uncovered/foreign activity before a global
elasticity is calculated. No empirical point estimate, universal sign or
complete solution of the research agenda is asserted.

\begin{thebibliography}{9}
\bibitem{DH} N. D\"obbeling-Hildebrandt and coauthors (2024), Systematic
review and meta-analysis of ex-post evaluations on the effectiveness of
carbon pricing, \emph{Nature Communications} 15, 4147.
\url{https://doi.org/10.1038/s41467-024-48512-w}.
Inspected publication PDF, printed p.\ 2 (price-level evidence) and
pp.\ 7--8 (discussion and policy-price distinctions):
\url{https://thedocs.worldbank.org/en/doc/83e2d33f048039a281426fa3bb150895-0050022025/original/Systematic-review-and-meta-analysis-of-ex-post-evaluations-on-the-effectiveness-of-carbon-pricing.pdf}.
\end{thebibliography}

\end{document}
